\chapter{紫外发散与超对称的解决之道}

\lectureoverview{在引入超对称的形式化数学之前，我们应当明确它旨在解决的是什么问题。我们将分别在位置空间和动量空间中描述紫外发散，通过量纲分析推导标量与费米子传播子的短距离行为，并估计对标量质量的主要圈修正。标量圈与费米子圈具有相同的二次截断依赖性，但符号相反。当我们追问究竟什么原理能使它们的耦合、多重数和质量匹配得足够好，以致这种抵消是结构性的而非偶然的，超对称才变得令人信服。}

\section{为什么要回到发散问题？}

本意是立即开始讲授超对称的机制。然而，最好还是将抽象推迟一讲，重新建立动机。先前关于发散和质量重整化的讲述进行得太快，除非我们首先确切地看到需要修补的是什么，否则超对称会显得像一项任意的数学发明。

费曼图以两种密切相关的方式进入粒子物理学。最为人熟知的是，它们组织散射振幅。我们指定入射和出射粒子，并对所需阶数的所有图所代表的振幅求和。它们也组织\emph{有效拉格朗日量}的构造。这里的``有效''一词并非表示异常成功的褒义词。它意味着这种描述是近似的：无法触及的短距离细节已被消去或积分掉，它们的影响被编码在剩余的参数和相互作用之中。

每一种这样的描述都有其适用范围。我们可以用一个最短的可信距离来指定这一范围，此处记作\(\delta\)，或者用一个最大的可信动量，
\begin{equation}
\Lambda\sim \delta^{-1}.
\end{equation}
这两个截断并非各自独立的假设。它们是同一紫外极限的傅里叶对偶描述。

\section{一个传播子，两种描述}

传播子是在一个时空点产生或插入一个粒子、并在另一个时空点将其移除或探测到的振幅。设这两个点为\(x^\mu\)和\(y^\mu\)，并将它们的间隔写作
\begin{equation}
\Delta^\mu=y^\mu-x^\mu.
\end{equation}
我们用小写符号\(\delta\)专指截断，这样间隔与正规化子就不会混淆。

平移不变性意味着自由传播子仅通过\(\Delta\)依赖于\(x\)和\(y\)。同一对象可以进行傅里叶变换：
\begin{equation}
G(\Delta)
=\int\frac{d^4k}{(2\pi)^4}\,e^{-ik\cdot\Delta}\,\widetilde G(k).
\end{equation}

当向听众提出``位置的傅里叶共轭是什么？''这一问题时，答案是动量。因此，将圈描述为回到起点的传播子这一位置空间描述，等价于将其描述为对绕圈循环的每一个动量进行积分这一动量空间描述。动量空间规则要求在每一个顶点施加动量守恒；位置空间则使当前的幂次计数论证更为透明。

在我们进行任何计算之前，这种等价性就已经指出了紫外问题。在位置空间中，圈探测的是重合点极限\(\Delta\to0\)。在动量空间中，它探测的是极限\(|k|\to\infty\)。大动量意味着短波长，因此短距离奇异性与高动量发散是同一现象在不同坐标下的表现。

本讲处理最简单的自旋类别：标量场、费米子场，以及作为背景的规范场。标量自旋为零，由一个单分量场\(\phi(x)\)表示。当本讲录制于2010年时，希格斯粒子尚未被发现，被描述为粒子物理学中预期的主要标量。这一历史背景解释了``希格斯-威格斯''这一戏称；该场在物理上极为迫切，但尚未在实验上得到确立。

对于实标量场，传播子可以示意性地写作
\begin{equation}
G_\phi(x,y)=\expect{0|T\!\left\{\phi(x)\phi(y)\right\}|0}
=G_\phi(\Delta).
\end{equation}
时间排序约定和归一化对我们的幂次计数无关紧要。重要的事实是，两个场插入分别描述在一个点处的产生和在另一个点处的湮灭。

\begin{classroomqa}
\audiencequestion{\(x\)和\(y\)是空间中的位置还是时空中的位置？\lecturetimestamp{00:09:31}}
\lecturerresponse{它们是时空点。每一个都是四维矢量。因此传播子依赖于它们的四维间隔，并且在洛伦兹不变的真空态中，依赖于相应的不变间隔。}
\end{classroomqa}

自旋为二分之一的粒子由狄拉克场\(\psi_i(x)\)描述。指标\(i\)是旋量指标，而不是空间分量，因此传播子是矩阵：
\begin{equation}
S_{ij}(x,y)=\expect{0|T\!\left\{\psi_i(x)\bar\psi_j(y)\right\}|0}
=S_{ij}(\Delta).
\end{equation}

向听众提出的一个问题给出了分量数：四维狄拉克场有四个复分量。无质量外尔场可以用两个复分量描述，因此左手场和右手场可以分别处理。就目前的量纲论证而言，重要的只是旋量矩阵结构的存在。

由于\(S_{ij}\)是\(4\times4\)矩阵，它可以按狄拉克矩阵的基展开。我们暂且把这些指标隐去，在需要时再恢复。这并不是声称矩阵结构不存在；而是在我们建立紫外幂律时有意作出的简化。

\begin{classroomqa}
\audiencequestion{共轭场或带匕首的场何时必须出现在传播子中？\lecturetimestamp{00:14:06}}
\lecturerresponse{在这个示意性讨论中，对于实的电中性标量场，同一个场既包含产生算符也包含湮灭算符，因此不需要一个不同的带匕首的场。对于带电场，必须区分粒子算符和反粒子算符。因此，对于电子，\(\psi\)、\(\psi^\dagger\)以及协变形式的\(\bar\psi=\psi^\dagger\gamma^0\)扮演着不同的角色。}
\end{classroomqa}

物理的希格斯粒子是电中性的，尽管底层的标准模型希格斯场是电弱二重态。讲座把实的中性标量场用作其局部模型；这一简化对于紫外估计已经足够。\editorialnote{这一区分防止把示意性的实标量例子误当作关于电弱希格斯多重态的完整陈述。}

\section{自然单位与标量场的量纲}

令\(\hbar=c=1\)。于是长度和时间具有相同的单位，而能量、质量和动量具有长度倒数的量纲：
\begin{equation}
[m]=[E]=[p]=L^{-1}.
\end{equation}
作用量出现在相位\(e^{iS}\)中，因此它是无量纲的。四维时空中一个有代表性的标量作用量为
\begin{equation}
S_\phi=\int d^4x\left[
\frac12\partial_\mu\phi\,\partial^\mu\phi
-\frac12m^2\phi^2
-\frac{\lambda}{4!}\phi^4
+\cdots\right].
\end{equation}
\(1/2\)和\(1/4!\)这些常规因子不影响量纲分析。

从动能项开始。由于\([d^4x]=L^4\)和\([\partial_\mu]=L^{-1}\)，无量纲性给出
\begin{align}
1
&=L^4\left(L^{-1}[\phi]\right)^2,\\
[\phi]&=L^{-1}.
\end{align}
质量项提供了一个检验：
\begin{equation}
L^4[m^2][\phi^2]
=L^4L^{-2}L^{-2}=1.
\end{equation}

接下来向听众提出的问题是四次耦合的量纲。它在四维时空中是无量纲的：
\begin{equation}
L^4[\lambda][\phi]^4=L^4[\lambda]L^{-4}=1
\quad\Longrightarrow\quad [\lambda]=L^0.
\end{equation}
这使得\(\phi^4\)相互作用在四维中很特殊。\(\phi^6\)及更高次相互作用的量纲遵循同样的规则，但它们的耦合带有质量的负幂，因此表明它们扮演着更受限制的有效理论角色。

因此，场的量纲已经告诉我们哪些相互作用可以以无量纲系数出现。同样的方法将确定传播子和汤川耦合，而无需计算详细的费曼积分。

\section{短距离处的标量传播子}

首先考虑无质量标量。不存在内禀长度标度，因此唯一可用的标度是间隔\(\Delta\)。平移不变性只留下\(\Delta\)，洛伦兹不变性把依赖关系约化为其不变模长，而两个标量场之积具有量纲\(L^{-2}\)。不计一个数值常数，
\begin{equation}
G_0(\Delta)\sim\frac{1}{\Delta^2}.
\end{equation}
诸如\(4\pi^2\)这样的因子无法仅由量纲推理确定。

质量引入康普顿长度\(m^{-1}\)。它的效应最容易从色散关系理解
\begin{equation}
E=\sqrt{\bm p^{,2}+m^2},
\qquad
E=|\bm p|\quad(m=0).
\end{equation}
在\(|\bm p|\gg m\)时，质量可以忽略。由于高动量探测短距离，有质量传播子必须在\(\Delta\to0\)时趋近无质量的\(1/\Delta^2\)行为。在远大于\(m^{-1}\)的距离上，传播受到压低。讲座用下式示意性地表示这一事实：
\begin{equation}
G_m(\Delta)\sim \frac{e^{-m|\Delta|}}{\Delta^2}.
\end{equation}
这有意不是一个精确公式。在四维欧几里得空间中，精确的自由传播子正比于\(mK_1(mr)/r\)，其中\(r=|\Delta|\)；它具有相同的\(1/r^2\)紫外奇异性，以及伴随幂律前因子的指数式长距离衰减。\editorialnote{贝塞尔函数形式为讲座的示意性指数提供了标准的精确补充，而讲座的幂次计数论证并不需要它。}

关键的教训很简单：质量软化长距离传播，但并不能修复重合点处的奇异性。在离一个点足够近时，任何有限质量都变得无关紧要。

\section{标量蝌蚪图与精细调节}

四次相互作用\(\lambda\phi^4\)产生一个四线顶点。把其中两条腿相互连接，留下两条外部标量腿。由此得到的蝌蚪图修正\(\phi^2\)的系数，从而修正标量质量的平方。

保留下来的画面捕捉到了从量纲分析到这一圈图修正的确切过渡。色散关系仍留在左侧，标量作用量和场量纲在右侧仍部分可见，而那个小蝌蚪图刚刚出现在中央。

\begin{figure}[H]
\centering
\includegraphics[width=0.72\textwidth]{lecture_03_figure_02.png}
\lectureframe{00:28:30.560}

\medskip
\begin{tikzpicture}[line width=0.8pt,scale=1.0]
  \draw[dashed] (-2.6,0) -- (0,0) -- (2.6,0);
  \fill (0,0) circle (0.055);
  \draw (0,0) .. controls (-0.85,0.7) and (-0.85,1.6) .. (0,1.6)
                .. controls (0.85,1.6) and (0.85,0.7) .. (0,0);
\end{tikzpicture}
\caption{标量自能过渡处的黑板，随后是对蝌蚪图拓扑的保守示意性重绘。截图可靠地显示了圈图、相对论性与无质量能量公式以及部分作用量；被遮挡的系数和标签不声称可见。}
\end{figure}

在顶点处，标量可以被吸收并重新发射，而另外两条腿则从该点传播回自身。除数值因子和组合因子外，修正为
\begin{equation}
\delta m_\phi^2\sim \lambda G_\phi(0).
\end{equation}
在连续近似下，重合的传播子是发散的。有效理论并不声称在任意短距离上有效，因此施加\(|\Delta|\geq\delta\)。于是
\begin{equation}
G_\phi(0)\longrightarrow G_\phi(\delta)sim\frac{1}{\delta^2},
\qquad
\delta m_\phi^2\sim\frac{\lambda}{\delta^2}sim\lambda\Lambda^2.
\end{equation}
该发散被称为二次发散，因为它包含动量截断的两个幂次。动量空间积分给出同样的结论：过大的贡献来自任意大的圈动量。最短尺度最终是由大统一、量子引力还是某种尚未想象到的东西设定，并不改变有效理论的诊断。

这是一个精细调节问题。如果\(\delta\)极其小，而\(\lambda\)并非异常微小，那么辐射修正就远大于低能下观测到的标量质量平方。裸系数与圈图修正随后必须抵消到令人难以置信的许多位数字：
\begin{equation}
m_{\mathrm{physical}}^2
=m_{\mathrm{bare}}^2+\delta m_\phi^2.
\end{equation}
有质量传播子并不能挽救局面。在导致发散的距离上，\(e^{-m|\Delta|}\simeq1\)。紫外圈图在到达截断之前早已忘记了质量。

\section{费米子量纲及其传播子}

我们现在对狄拉克费米子重复同样的量纲分析。其自由作用量为
\begin{equation}
S_\psi=\int d^4x\,
\bar\psi\left(i\gamma^\mu\partial_\mu-m\right)\psi.
\end{equation}
伽马矩阵是数值的且无量纲。动能项只包含一个导数，因此
\begin{align}
1
&=L^4[\bar\psi]L^{-1}[\psi],\\
[\psi]=[\bar\psi]&=L^{-3/2}.
\end{align}
半整数幂不应令人惊慌；费米子把半整数带到了许多地方。

\begin{classroomqa}
\audiencequestion{从质量项推导\([\psi]\)不是更容易吗？\lecturetimestamp{00:36:02}}
\lecturerresponse{是的。质量项给出\(L^4L^{-1}[\bar\psi\psi]=1\)，因此\([\bar\psi\psi]=L^{-3}\)以及\([\psi]=L^{-3/2}\)。之所以先用动能项推导，是因为它在无质量理论中仍然有意义，并显示出为什么费米子质量表现为\(m\)，而不是\(m^2\)。}
\end{classroomqa}

两个费米子场的量纲为\(L^{-3}\)，因此讲义中对无质量传播子的紧凑量纲简写为
\begin{equation}
S(\Delta)\sim\frac{1}{|\Delta|^3}.
\end{equation}
恢复其矩阵结构后，就得到洛伦兹协变性所要求的更精确的分子：
\begin{equation}
S_{ij}(\Delta)
\propto
\frac{\left(\gamma_\mu\Delta^\mu\right)_{ij}}
{\left(\Delta^2\right)^2}.
\end{equation}
分子增加了一个长度幂次，因此分母包含间隔的四个幂次；净量纲仍为\(L^{-3}\)。与标量情形一样，费米子质量在大间隔处压低传播子，但不改变这种短距离标度行为。

\section{为什么闭合费米子圈是负的}

我们现在可以问，费米子图是否可能抵消糟糕的玻色子自能。标量蝌蚪图具有正号，至多相差诸如\(\pi\)的幂次和组合系数这样的正数值因子。闭合费米子圈则带有相反的符号。

\begin{classroomqa}
\audiencequestion{费米子圈的负号与先前关于一整圈旋转的讨论有关吗？\lecturetimestamp{00:40:21}}
\lecturerresponse{它属于同一费米子符号结构，但这里有用的证明是通过交换反对称性，而不是通过圈图的几何取向。切开两个圈，交换全同粒子，再把它们重新连接起来。交换把两个圈变成一个圈，并且只对费米子贡献一个负号。}
\end{classroomqa}

这个论证值得完整给出。假设某个单圈图具有符号\(\sigma\)，即正或负。两个不相连的全同圈具有符号
\begin{equation}
\sigma^2=+1.
\end{equation}
同时切开两个圈。我们现在有两个全同粒子在切口之间传播。它们通常的两粒子波函数对玻色子是交换对称的，对费米子是交换反对称的：
\begin{align}
\Psi_B(x_1,x_2)&=+\Psi_B(x_2,x_1),\\
\Psi_F(x_1,x_2)&=-\Psi_F(x_2,x_1).
\end{align}
反对称性也解释了泡利原理：如果两个费米子被置于同一状态，交换在物理上不改变任何东西，但必须反转波函数，从而迫使它为零。

交换这两个粒子后，将切断的线重新连接。从拓扑上看，两个圈已经变成了一个。未交换的双圈振幅为正。交换使玻色子的振幅保持为正，而使费米子的振幅反号。因此
\begin{equation}
\begin{aligned}
\operatorname{sign}(\text{closed boson loop})&=+1,\\
\operatorname{sign}(\text{closed fermion loop})&=-1.
\end{aligned}
\end{equation}
这个结论是普遍的：每个独立的闭合费米子圈贡献一个因子\(-1\)。

\section{汤川圈}

回到标量质量。一个玻色子不能转变为单个费米子，因为费米子数宇称会改变。最简单的费米子修正中，标量分裂为一个费米子-反费米子对，该对绕圈传播并重新结合。这种相互作用是汤川耦合，
\begin{equation}
\mathcal L_Y=-g\phi\bar\psi\psi.
\end{equation}
对于希格斯，费米子可以是夸克或轻子。有两个相互作用顶点，因此有两个\(g\)的因子。

标量蝌蚪图涉及一个顶点，其整体位置可以固定，因为平移不变性仅仅提供时空体积。汤川圈涉及两个顶点。我们可以固定一个顶点，或者等价地固定它们的平均位置，但我们必须对它们的相对间距\(\Delta\)积分。每个费米子传播子贡献短距离因子\(1/|\Delta|^3\)，所以
\begin{equation}
\delta m_{\phi,F}^2
\sim
-g^2\int_{|\Delta|\geq\delta}d^4\Delta\,
\frac{1}{|\Delta|^6}.
\end{equation}

在两个传播子、间距积分和两个耦合都写出之后，还差最后一个因子。答案是负号。两个传播子闭合成一个费米子圈，因此整个原本为正的短距离积分要乘以\(-1\)。

量纲分析已经确定了发散。四个长度幂来自\(d^4\Delta\)，分母中出现六个，而唯一剩下的长度标度是截断：
\begin{equation}
\int_{|\Delta|\geq\delta}
\frac{d^4\Delta}{|\Delta|^6}
\sim\frac{1}{\delta^2}.
\end{equation}
在四维欧几里得径向坐标中，这可以明确写出：
\begin{equation}
\int_{\delta}^{R}d^4\Delta\,\frac{1}{|\Delta|^6}
=2\pi^2\int_{\delta}^{R}dr\,\frac{r^3}{r^6}
=\pi^2\left(\frac{1}{\delta^2}-\frac{1}{R^2}\right).
\end{equation}
旋量迹和传播子归一化会改变系数，但不会改变截断的二次幂。\editorialnote{径向积分是讲座中量纲估计的一个标准显式版本。它用欧几里得距离写出，在那里正定性和紫外幂次计数是透明的。}

\begin{classroomqa}
\audiencequestion{反转或选择积分限会改变符号吗？\lecturetimestamp{00:51:06}}
\lecturerresponse{不会。这不是一个由积分限产生物理符号的一维有向积分。它是一个对四个分量的积分，或者更干净地说，是对一个正径向测度的积分。它的紫外大小为正。负的贡献完全来自闭合费米子圈。}
\end{classroomqa}

\begin{classroomqa}
\audiencequestion{汤川耦合\(g\)的量纲是什么？\lecturetimestamp{00:51:51}}
\lecturerresponse{它在四维中是无量纲的。由于\([\phi]=L^{-1}\)和\([\bar\psi\psi]=L^{-3}\)，乘积\(\phi\bar\psi\psi\)具有量纲\(L^{-4}\)，恰好被\(d^4x\)抵消。}
\end{classroomqa}

因此，领头费米子修正为
\begin{equation}
\delta m_{\phi,F}^2
\sim -C_F\frac{g^2}{\delta^2}
=-C_Fg^2\Lambda^2,
\end{equation}
其中\(C_F>0\)包含数值、自旋和多重度因子。标量圈具有相应的形式
\begin{equation}
\delta m_{\phi,B}^2
\sim+C_B\frac{\lambda}{\delta^2}
=+C_B\lambda\Lambda^2.
\end{equation}
我们还没有解决精细调节问题，但我们已经发现了某种不可或缺的东西：一个具有相同紫外量级且符号相反的贡献。

\section{抵消所需的条件}

如果无关的耦合碰巧满足
\begin{equation}
C_B\lambda=C_Fg^2
\end{equation}
且精确到大的标度等级所要求的精度，那将是一个非凡的巧合。讲座将其简写为\(g^2=\lambda\)，省略了依赖于约定和多重度的因子。真正的解决方案需要一个数学结构来强制这一关系，而不仅仅是允许它。

这就是超对称的动机。超对称将玻色子和费米子联系起来，并限制它们的耦合为使得许多原本危险的发散项相消的组合。然后这种抵消会在各个图之间重复出现，因为它源于一种对称性，而不是因为参数被一个圈一个圈地调整。

质量匹配同样重要。如果玻色子和费米子的质量完全相同，它们的完整传播子可以最紧密地配对。如果质量略有不同，完整的有限答案不会精确抵消，因为传播子在与其康普顿波长相当的距离上存在差异。然而，它们共同的短距离渐近行为仍然可以抵消发散部分：
\begin{equation}
\mathcal I_B(\Delta)-\mathcal I_F(\Delta)
\quad\hbox{may be finite as}\quad \Delta\to0,
\end{equation}
其中\(\mathcal I_B\)和\(\mathcal I_F\)表示在包含顶点、分子和多重数之后完整的玻色子和费米子圈被积函数。传播子本身具有不同的量纲，并不是被直接减除的。即使紫外项抵消了，也可能残留一个有限的剩余项。这正是近似超对称或破缺超对称的切入点。精确相等给出最干净的抵消；足够受控的破缺仍可能消除紫外灾难，同时留下可观测的质量劈裂。

因此，我们还没有写出超对称的代数，就已经到达了超对称。推理链条是完整的：
\begin{equation}
\begin{gathered}
G_\phi(\Delta)\sim \Delta^{-2},
\qquad
\delta m_{\phi,B}^2\sim+C_B\lambda\Lambda^2,\\[0.4em]
S_\psi(\Delta)\sim \Delta^{-3},
\qquad
\delta m_{\phi,F}^2\sim-C_Fg^2\Lambda^2,\\[0.4em]
\text{symmetry-enforced matching}\\
\Longrightarrow\quad
\text{quadratic divergences cancel}.
\end{gathered}
\end{equation}
接下来的任务是理解强制这些关系的抽象对称原理，并看到它如何作用的具体例子。